% GENERATED FILE. Edit canonical lessons, then run npm run generate:books.
% canonical-source-hash: fnv1a64:dbd67942275bcbc4
% canonical-lessons: JA-C424-otoshimono, JA-C424-hagureru, JA-C424-miushinau, JA-C424-sagashimawaru, JA-C424-saiwai

\chapter{Lost and Found}
\label{ch:ja-lost-and-found}
\hlchaptermodality{424}

\begin{chapteropening}
\textbf{By the end of this chapter:} \emph{I can talk about losing people and things on a trip.}

Complete the last lesson of chapter 424: I can talk about losing people and things on a trip.
\end{chapteropening}

\section[\texorpdfstring{\ja{おとしもの} (otoshimono)}{おとしもの (otoshimono)}]{\ja{おとしもの} (otoshimono) — a lost item; something dropped}

\label{lesson:JA-C424-otoshimono}



\begin{quote}
Before the new one: say the Japanese for to treat someone, then the Japanese for a full stomach.
\end{quote}

\subsection*{You'll want to know: \ja{おとしもの}}
\textbf{\ja{おとしもの}} — \emph{otoshimono} — \textquotedblleft{}a lost item; something dropped\textquotedblright{}.

\begin{quote}
\textbf{\ja{こうばんに} \ja{おとしものを} \ja{とどけました。}} — \emph{kōban ni otoshimono o todokemashita} — \textquotedblleft{}I took something I had found to the police box.\textquotedblright{}
\end{quote}

\begin{cousinweb}
\textbf{\ja{おとしもの}} (\emph{otoshimono}) is a native Japanese word (wago), inherited rather than borrowed from Chinese, built from \emph{otoshi}, \textquotedblleft{}dropping\textquotedblright{} (from \emph{otosu}) and \emph{mono}, \textquotedblleft{}thing\textquotedblright{}; its kanji were attached later as kun'yomi (JMdict).
\end{cousinweb}

\begin{grammarlens}[title={What you've built}]
One more word for getting around.
\end{grammarlens}

\subsection*{Your turn}
\begin{itemize}
\raggedright
  \item \emph{Say it:} \emph{otoshimono}
  \item \emph{Say it:} \emph{otoshimono}, once more
  \item \emph{Say it:} \emph{kōban ni otoshimono o todokemashita} — I took something I had found to the police box.
  \item \emph{Recall:} say the Japanese for to treat someone, then the Japanese for a full stomach, then say \emph{otoshimono} again
\end{itemize}

\begin{tcolorbox}[breakable,colback=teal!4,colframe=teal!35!black,title={Before you move on}]
What does \ja{おとしもの} mean? (\textquotedblleft{}A lost item; something dropped\textquotedblright{}.) Say it once more.
\end{tcolorbox}
\section[\texorpdfstring{\ja{はぐれる} (hagureru)}{はぐれる (hagureru)}]{\ja{はぐれる} (hagureru) — to get separated (from one's group)}

\label{lesson:JA-C424-hagureru}



\begin{quote}
Before the new one: say the Japanese for a full stomach, then the Japanese for a lost item.
\end{quote}

\subsection*{You'll want to know: \ja{はぐれる}}
\textbf{\ja{はぐれる}} — \emph{hagureru} — \textquotedblleft{}to get separated (from one's group)\textquotedblright{}.

\begin{quote}
\textbf{\ja{まつりで} \ja{かぞくと} \ja{はぐれました。}} — \emph{matsuri de kazoku to haguremashita} — \textquotedblleft{}I got separated from my family at the festival.\textquotedblright{}
\end{quote}

Polite \textbf{\ja{はぐれます}}, \emph{haguremasu}; an ichidan verb. Past (polite): \textbf{\ja{こどもが} \ja{ははと} \ja{はぐれました。}}, \emph{kodomo ga haha to haguremashita}.

\begin{cousinweb}
\textbf{\ja{はぐれる}} (\emph{hagureru}) is a native Japanese word (wago), inherited rather than borrowed from Chinese (JMdict).
\end{cousinweb}

\begin{grammarlens}[title={What you've built}]
One more word for getting around.
\end{grammarlens}

\subsection*{Your turn}
\begin{itemize}
\raggedright
  \item \emph{Say it:} \emph{hagureru}
  \item \emph{Say it:} \emph{hagureru}, once more
  \item \emph{Say it:} \emph{matsuri de kazoku to haguremashita} — I got separated from my family at the festival.
  \item \emph{Recall:} say the Japanese for a full stomach, then the Japanese for a lost item, then say \emph{hagureru} again
\end{itemize}

\begin{tcolorbox}[breakable,colback=teal!4,colframe=teal!35!black,title={Before you move on}]
What does \ja{はぐれる} mean? (\textquotedblleft{}To get separated (from one's group)\textquotedblright{}.) Say it once more.
\end{tcolorbox}
\section[\texorpdfstring{\ja{みうしなう} (miushinau)}{みうしなう (miushinau)}]{\ja{みうしなう} (miushinau) — to lose sight of}

\label{lesson:JA-C424-miushinau}



\begin{quote}
Before the new one: say the Japanese for a lost item, then the Japanese for to get separated.
\end{quote}

\subsection*{You'll want to know: \ja{みうしなう}}
\textbf{\ja{みうしなう}} — \emph{miushinau} — \textquotedblleft{}to lose sight of\textquotedblright{}.

\begin{quote}
\textbf{\ja{えきで} \ja{ともだちを} \ja{みうしないました。}} — \emph{eki de tomodachi o miushinaimashita} — \textquotedblleft{}I lost sight of my friend at the station.\textquotedblright{}
\end{quote}

Polite \textbf{\ja{みうしないます}}, \emph{miushinaimasu}; a godan verb. Past (polite): \textbf{\ja{ははが} \ja{こどもを} \ja{みうしないました。}}, \emph{haha ga kodomo o miushinaimashita}.

\begin{cousinweb}
\textbf{\ja{みうしなう}} (\emph{miushinau}) is a native Japanese word (wago), inherited rather than borrowed from Chinese, built from \emph{mi}, \textquotedblleft{}seeing\textquotedblright{} (from \emph{miru}) and \emph{ushinau}, \textquotedblleft{}to lose\textquotedblright{}; its kanji were attached later as kun'yomi (JMdict).
\end{cousinweb}

\begin{grammarlens}[title={What you've built}]
One more word for getting around.
\end{grammarlens}

\subsection*{Your turn}
\begin{itemize}
\raggedright
  \item \emph{Say it:} \emph{miushinau}
  \item \emph{Say it:} \emph{miushinau}, once more
  \item \emph{Say it:} \emph{eki de tomodachi o miushinaimashita} — I lost sight of my friend at the station.
  \item \emph{Recall:} say the Japanese for a lost item, then the Japanese for to get separated, then say \emph{miushinau} again
\end{itemize}

\begin{tcolorbox}[breakable,colback=teal!4,colframe=teal!35!black,title={Before you move on}]
What does \ja{みうしなう} mean? (\textquotedblleft{}To lose sight of\textquotedblright{}.) Say it once more.
\end{tcolorbox}
\section[\texorpdfstring{\ja{さがしまわる} (sagashimawaru)}{さがしまわる (sagashimawaru)}]{\ja{さがしまわる} (sagashimawaru) — to search all over}

\label{lesson:JA-C424-sagashimawaru}



\begin{quote}
Before the new one: say the Japanese for to get separated, then the Japanese for to lose sight of.
\end{quote}

\subsection*{You'll want to know: \ja{さがしまわる}}
\textbf{\ja{さがしまわる}} — \emph{sagashimawaru} — \textquotedblleft{}to search all over\textquotedblright{}.

\begin{quote}
\textbf{\ja{さいふを} \ja{さがしまわりました。}} — \emph{saifu o sagashimawarimashita} — \textquotedblleft{}I searched everywhere for my wallet.\textquotedblright{}
\end{quote}

Polite \textbf{\ja{さがしまわります}}, \emph{sagashimawarimasu}; a godan verb. Past (polite): \textbf{\ja{あねが} \ja{かぎを} \ja{さがしまわりました。}}, \emph{ane ga kagi o sagashimawarimashita}.

\begin{cousinweb}
\textbf{\ja{さがしまわる}} (\emph{sagashimawaru}) is a native Japanese word (wago), inherited rather than borrowed from Chinese, built from \emph{sagashi}, \textquotedblleft{}searching\textquotedblright{} (from \emph{sagasu}) and \emph{mawaru}, \textquotedblleft{}to go round\textquotedblright{}; its kanji were attached later as kun'yomi (JMdict).
\end{cousinweb}

\begin{grammarlens}[title={What you've built}]
One more word for getting around.
\end{grammarlens}

\subsection*{Your turn}
\begin{itemize}
\raggedright
  \item \emph{Say it:} \emph{sagashimawaru}
  \item \emph{Say it:} \emph{sagashimawaru}, once more
  \item \emph{Say it:} \emph{saifu o sagashimawarimashita} — I searched everywhere for my wallet.
  \item \emph{Recall:} say the Japanese for to get separated, then the Japanese for to lose sight of, then say \emph{sagashimawaru} again
\end{itemize}

\begin{tcolorbox}[breakable,colback=teal!4,colframe=teal!35!black,title={Before you move on}]
What does \ja{さがしまわる} mean? (\textquotedblleft{}To search all over\textquotedblright{}.) Say it once more.
\end{tcolorbox}
\section[\texorpdfstring{\ja{さいわい} (saiwai)}{さいわい (saiwai)}]{\ja{さいわい} (saiwai) — fortunately; luckily}

\label{lesson:JA-C424-saiwai}



\begin{quote}
Before the new one: say the Japanese for to lose sight of, then the Japanese for to search all over.
\end{quote}

\subsection*{You'll want to know: \ja{さいわい}}
\textbf{\ja{さいわい}} — \emph{saiwai} — \textquotedblleft{}fortunately; luckily\textquotedblright{}.

\begin{quote}
\textbf{\ja{さいわい、さいふは} \ja{こうばんに} \ja{ありました。}} — \emph{saiwai, saifu wa kōban ni arimashita} — \textquotedblleft{}Luckily, my wallet was at the police box.\textquotedblright{}
\end{quote}

\begin{cousinweb}
\textbf{\ja{さいわい}} (\emph{saiwai}) is a native Japanese word (wago), inherited rather than borrowed from Chinese; its kanji was attached later as a kun'yomi (JMdict).
\end{cousinweb}

\begin{grammarlens}[title={What you've built}]
One more word for getting around.
\end{grammarlens}

\subsection*{Your turn}
\begin{itemize}
\raggedright
  \item \emph{Say it:} \emph{saiwai}
  \item \emph{Say it:} \emph{saiwai}, once more
  \item \emph{Say it:} \emph{saiwai, saifu wa kōban ni arimashita} — Luckily, my wallet was at the police box.
  \item \emph{Recall:} say the Japanese for to lose sight of, then the Japanese for to search all over, then say \emph{saiwai} again
\end{itemize}

\begin{tcolorbox}[breakable,colback=teal!4,colframe=teal!35!black,title={Before you move on}]
What does \ja{さいわい} mean? (\textquotedblleft{}Fortunately; luckily\textquotedblright{}.) Say it once more.
\end{tcolorbox}
